% SOURCE_STATEMENT: "For a globally hyperbolic expanding vacuum spacetime with compact Cauchy hypersurfaces and a CMC foliation, what is the asymptotic behavior as the past singularity is approached, and what is the future behavior of the spacetime?"
% Preparation and final self-audit: CONTRIBUTING.md prompts C and D; source/audit artifacts retained in the local physics_statements_2026_10_11 batch.
% !TeX program = lualatex
\documentclass[11pt]{article}
\usepackage[margin=25mm]{geometry}
\usepackage{fontspec}
\setmainfont{Latin Modern Roman}
\usepackage{amsmath,amssymb,mathtools}
\usepackage{unicode-math}
\setmathfont{Latin Modern Math}
\usepackage{xurl,hyperref}
\hypersetup{colorlinks=true,urlcolor=blue}
\setcounter{secnumdepth}{0}
\setlength{\parindent}{0pt}
\setlength{\parskip}{5pt}
\setlength{\emergencystretch}{3em}
\newcommand{\sref}[2]{\hyperlink{source:#1:#2}{[S#2]}}
\newcommand{\cataloguescope}[1]{\paragraph{Recorded target.}#1}
\newcommand{\R}{\mathbb R}
\newcommand{\C}{\mathbb C}
\newcommand{\N}{\mathbb N}
\newcommand{\Z}{\mathbb Z}
\newcommand{\Q}{\mathbb Q}
\newcommand{\F}{\mathbb F}
\newcommand{\eps}{\varepsilon}
\begin{document}
% UnsolvedMath ID 30004803; source catalogue code OWR-8415343-006
\setcounter{section}{30004802}
\section{Asymptotics of Expanding Vacuum Spacetimes Near Singularities}\label{30004803}
{\small UnsolvedMath ID 30004803 · Mathematical Physics}
\subsection{Definitions and mathematical statement}

\paragraph{Shared notation.}$\mathbb N=\{1,2,\ldots\}$; the standard number fields have their usual meanings. All additional conventions are stated locally.


\paragraph{Definitions and conventions.}
A globally hyperbolic spacetime has Cauchy hypersurfaces crossed once by every inextendible causal curve. Vacuum means zero matter stress tensor. A constant-mean-curvature (CMC) foliation uses spatial slices on which the trace of their second fundamental form is spatially constant; its sign convention and monotone time parameter describe expansion. Compact Cauchy slices describe a spatially closed cosmological model. Past-singularity asymptotics and future expanding behavior are distinct temporal limits.

\paragraph{Target retained from the canonical source catalogue.}
\paragraph{Catalogue wording retained for traceability.}
For a globally hyperbolic expanding vacuum spacetime with compact Cauchy hypersurfaces and a CMC foliation, what is the asymptotic behavior as the past singularity is approached, and what is the future behavior of the spacetime?

\paragraph{Source scope and qualification.}
The report discusses the $3+1$-dimensional setting, with some results extending further. No particular rescaled limit, topology or genericity condition is fixed by the extraction; both temporal directions remain in the research target.
\subsection{Short English statement}
Describe both the approach to the past singularity and the future evolution of expanding vacuum cosmologies with compact CMC slices.
\subsection{Sources}
\begin{itemize}
\item[\hypertarget{source:30004803:1}{S1}] ULAM.AI and the UnsolvedMath Contributors, UnsolvedMath, record 30004803 (OWR-8415343-006). Catalogue curation under CC BY 4.0; definitions and English exposition added locally.
\url{https://huggingface.co/datasets/ulamai/UnsolvedMath}.
\item[\hypertarget{source:30004803:2}{S2}] Mathematical Aspects of General Relativity (2021), Oberwolfach Reports, DOI 10.4171/owr/2021/40, pp. 2157--2267. Primary report accessed 2026-10-11; relevant PDF pages 42.
\url{https://doi.org/10.4171/owr/2021/40}.
\item[\hypertarget{source:30004803:3}{S3}] Lott, initial geometry of vacuum cosmological spacetimes (catalogue-listed context; not independently consulted in this preparation).
\url{https://arxiv.org/abs/1908.02185}.
\item[\hypertarget{source:30004803:4}{S4}] Lott, future behavior of expanding spacetimes (catalogue-listed context; not independently consulted in this preparation).
\url{https://arxiv.org/abs/1709.03952}.
\end{itemize}
\label{end:30004803}
\end{document}
